\documentclass[10pt]{article}
\usepackage[T1]{fontenc}
\usepackage[utf8]{inputenc}
\usepackage[french]{babel}
\usepackage{amsmath,amssymb,mathrsfs,mathpazo}
\usepackage{geometry}
\geometry{a4paper,margin=2cm}
\usepackage{graphicx,array,tabularx,tabularray,colortbl,multirow,multicol,booktabs,enumitem,pifont,ulem,fancybox,diagbox,cancel}
\usepackage{pstricks,pstricks-add,pst-plot,pst-tree,pst-node,pst-text,pst-func,pst-math,pst-bspline,pst-3dplot,pst-eucl,pst-solides3d}
\let\pscalculate\undefined
\usepackage{pst-bezier}
\usepackage{tikz,pgfplots,tkz-tab}
\usetikzlibrary{arrows,arrows.meta,patterns,calc,positioning,trees,decorations.markings,shapes.geometric,shapes.misc,shapes.arrows,intersections,decorations.pathmorphing,decorations.pathreplacing,backgrounds,shadows,babel}
\pgfplotsset{compat=1.18}
\usepgfplotslibrary{fillbetween}
\usepackage[np]{numprint}
\usepackage{listings,makecell,pas-tableur,twoopt}
\let\barre\undefined
\usepackage[pstricks]{tablvar}
\definecolor{darkgreen}{RGB}{0,100,0}
\providecommand{\red}{\color{red}}
\providecommand{\blue}{\color{blue}}
\providecommand{\green}{\color{green}}
\providecommand{\black}{\color{black}}
\providecommand{\white}{\color{white}}
\providecommand{\gray}{\color{gray}}
\providecommand{\cyan}{\color{cyan}}
\providecommand{\magenta}{\color{magenta}}
\providecommand{\yellow}{\color{yellow}}
\providecommand{\violet}{\color{violet}}
\providecommand{\orange}{\color{orange}}
\providecommand{\euro}{\texteuro}
\providecommand{\EUR}{\texteuro}
\providecommand{\ieme}{\textsuperscript{e}}
\providecommand{\degres}{\ensuremath{^\circ}}
\providecommand{\degree}{\ensuremath{^\circ}}
\providecommand{\vv}[1]{\overrightarrow{#1}}
\providecommand{\Coord}[1]{#1}
\providecommand{\decosix}{}
\providecommand{\starredbullet}{\ensuremath{\bullet}}
\pagestyle{empty}
\setlength{\parindent}{0pt}
\usepackage{adjustbox}\DeclareUnicodeCharacter{2236}{:}\DeclareUnicodeCharacter{2190}{\ensuremath{\leftarrow}}
\providecommand{\R}{\mathbb{R}}
\providecommand{\N}{\mathbb{N}}
\providecommand{\Z}{\mathbb{Z}}
\providecommand{\Q}{\mathbb{Q}}
\providecommand{\D}{\mathbb{D}}
\providecommand{\C}{\mathbb{C}}
\providecommand{\e}{{\mathrm{e}}}
\providecommand{\vect}[1]{\overrightarrow{#1}}
\providecommand{\vectt}[1]{\overrightarrow{\mathrm{#1}}}
\providecommand{\vvt}[1]{\overrightarrow{\mathrm{#1}}}
\providecommand{\np}[2][]{{#2~#1}}
\providecommand{\vv}[1]{{\overrightarrow{#1}}}
\renewcommand{\d}{{\mathrm{d}}}
\providecommand{\degres}{{^\circ}}
\providecommand{\degre}{{^\circ}}
\providecommand{\blue}{\color{blue}}
\providecommand{\red}{\color{red}}
\providecommand{\enskip}{\quad}
\providecommand{\llbracket}{\lbrack\!\lbrack}
\providecommand{\rrbracket}{\rbrack\!\rbrack}
\providecommand{\euro}{{\text{€}}}
\providecommand{\no}{{\text{n°}}}
\providecommand{\interval}[2]{{[#1;#2]}}
\providecommand{\intervFO}[2]{{[#1;#2[}}
\providecommand{\dotfill}{\ldots}

\providecommand{\ds}{\displaystyle}
\providecommand{\textding}[1]{#1}
\providecommand{\Oij}{\(\mathrm{O};\vec{\imath},\vec{\jmath}\)}
\providecommand{\Oijk}{\(\mathrm{O};\vec{\imath},\vec{\jmath},\vec{k}\)}
\providecommand{\Ouv}{\(\mathrm{O};\vec{u},\vec{v}\)}
\providecommand{\intervFF}[2]{\ensuremath{[#1;#2]}}
\renewcommand{\euro}{\ifmmode\text{\texteuro}\else\texteuro\fi}
\DeclareUnicodeCharacter{00B0}{\ifmmode^{\circ}\else\textdegree\fi}
\DeclareUnicodeCharacter{20AC}{\ifmmode\text{\texteuro}\else\texteuro\fi}
% Énoncés : Éducation nationale. Transcription et corrigés : APMEP, Denis Vergès et contributeurs. Découpage et classement : CoopMaths, Sébastien Lozano, Éric Elter et contributeurs. Adaptation de la présentation : Maths Explorer.
% CC BY-SA : https://coopmaths.fr/www/about/
% https://coopmaths.fr/alea/static/bac/2023/tex/bac_2023_03_sujet2_polynesie_4_cor.tex
\begin{document}
\noindent\textit{APMEP / CoopMaths. CC BY-SA. Présentation adaptée par Maths Explorer.}\par\medskip
\medskip

\begin{enumerate}
\item Affirmation 1 : Vraie

Pour tout entier naturel $n$ : $n\geqslant 0 \Longrightarrow n+1 \geqslant 1 \Longrightarrow \dfrac{1}{n+1} \leqslant 1$  et  $-\dfrac{1}{n+1} \geqslant -1$.

De plus,  pour tout entier naturel $n$, $-1 \leqslant (-1)^n \leqslant 1 $.

D'où $ \dfrac{-1}{n+1} \leqslant \dfrac{(-1)^n}{n+1} \leqslant \dfrac{1}{n+1}$ car $n+1 >0$

finalement on a : $-1 \leqslant \dfrac{-1}{n+1} \leqslant \dfrac{(-1)^n}{n+1} \leqslant \dfrac{1}{n+1}\leqslant 1$.

La suite $u$ est bornée par -1 et 1.

\item Affirmation 2 : Fausse

Pour montrer qu'une affirmation est fausse, il suffit de donner un contre-exemple :

Considérons la suite $u$ définie sur $\N$ par $u_n= (-1)^n$.

Cette suite est bornée par $-1$ et 1 mais elle n'est pas convergente.

\item Affirmation 3 : Fausse

Pour montrer qu'une affirmation est fausse, il suffit de donner un contre-exemple :

Considérons la suite $u$ définie sur $\N$ par $u_n=-\dfrac{1}{n}$.

Cette suite est croissante et converge vers 0.

\item Affirmation 4 : Fausse 

Pour étudier la convexité d'une fonction,  il faut déterminer le signe de sa dérivée seconde. 

$f$ est dérivable sur $\R$ comme composée de fonctions dérivables.

$f$ est de la forme $\ln(u)$ avec $u(x)=x^2+2x+2$.  On a $u'(x)=2x+2$.

$f'=\dfrac{u'}{u}$ donc, pour tout réel $x$, $f'(x)=\dfrac{2x+2}{x^2+2x+2}$.

$f'$ est de la forme $\dfrac{v}{w}$ avec $v(x)=2x+2$ et $w(x)=x^2+2x+2$.

On a $v'(x)=2$ et $w'(x)=2x+2$.

$f''= \dfrac{v'\times w - w'\times v}{w^2}$  donc,  pour tout réel $x$,  $f''(x) =\dfrac { 2\times (x^2+2x+2)- (2x+2)\times(2x+2)    }{ {\left( x^2+2x+2 \right)}^2}$ 

$f''(x)=\dfrac {  2x^2+4x+4-4x^2-8x-4  }{ {\left( x^2+2x+2 \right)}^2} =\dfrac {  -2x^2-4x  }{ {\left( x^2+2x+2 \right)}^2} =\dfrac{-2x(x+2)}{ {\left( x^2+2x+2 \right)}^2}$.

Un carré est toujours strictement positif, $f''(x)$ est donc du signe de $-2x(x+2)$ qui est une fonction polynôme du second degré avec un coefficient dominant négatif (-2) et qui a pour racines 0 et -2.

On a donc le tableau de signe suivant :


\begin{center}
\begin{tikzpicture}[scale=0.875]

\tikzstyle{cadre}=[thin]
\tikzstyle{fleche}=[->,>=latex,thin]
\tikzstyle{nondefini}=[lightgray]

\def\Lrg{1.5}
\def\HtX{1}
\def\HtY{0.5}

\def\lignex{-0.5*\HtX}
\def\lignef{-1.5*\HtX}
\def\separateur{-0.5*\Lrg}

\def\gauche{-1.5*\Lrg}
\def\droite{6.5*\Lrg}

\def\haut{0.5*\HtX}
\def\bas{-2.5*\HtX-2*\HtY}

\node at (-1*\Lrg,0) {$x$};
\node at (0*\Lrg,0) {$-\infty$};
\node at (2*\Lrg,0) {$-2$};
\node at (4*\Lrg,0) {$0 $};
\node at (6*\Lrg,0) {$+\infty$};

\node at (-1*\Lrg,-1*\HtX) {$f''(x)$};
\node at (0*\Lrg,-1*\HtX) {$$};
\node at (1*\Lrg,-1*\HtX) {$-$};
\node at (2*\Lrg,-1*\HtX) {$0$};
\node at (3*\Lrg,-1*\HtX) {$+$};
\node at (4*\Lrg,-1*\HtX) {$0$};
\node at (5*\Lrg,-1*\HtX) {$-$};
\node at (6*\Lrg,-1*\HtX) {$$};


\draw[cadre] (\separateur,\haut) -- (\separateur, \lignef);
\draw[cadre] (\gauche,\haut) rectangle  (\droite, \lignef);
\draw[cadre] (\gauche,\lignex) -- (\droite,\lignex);
\end{tikzpicture}
\end{center}











On a déduit que $f$ est concave sur $]-\infty~;~-2]$ et sur $[0~;~+\infty[$,  et elle est convexe sur $[-2~;~0]$.

Elle n'est donc pas convexe sur $[-3~;~1]$.

\item  Affirmation 5 :  vraie

On initialise $M$ avec le premier terme de la suite, puis l'algorithme  va comparer chaque terme de la liste à $M$,  si un terme est supérieur à $M$, il va remplacer $M$.

Cet algorithme donne donc la valeur maximale des termes de la suite, soit 7 dans notre cas.
\end{enumerate}

\end{document}
